\documentclass[10pt,a4paper]{amsart}
\usepackage[margin=19mm]{geometry}
\usepackage{amsmath,amssymb,mathtools}
\usepackage[scr=rsfs,cal=euler]{mathalfa}
\usepackage{xfrac}
\usepackage[unicode,hidelinks,bookmarks=false]{hyperref}
\newcommand{\ie}{i.e.\ }
\newcommand{\cf}{cf.\ }
\newcommand{\resp}{respectively\ }
\newcommand{\Subsection}{\subsection}
\providecommand{\nobreakdash}{\nobreak\hbox{-}}
\setlength{\emergencystretch}{2em}
\multlinegap0pt
\usepackage{amssymb}
\usepackage{xcolor}
\newtheorem{proposition}{Proposition}
\def\l{\lambda}
\title{The Faber-Krahn inequality for $p$-Hermite operators}
\author{Xiaomei Sun, Kui Wang, Anqiang Zhu}
\date{Source: arXiv:2609.01269v1 (CC BY 4.0)}
\begin{document}
\maketitle

\noindent\emph{Source and licence.} The Faber-Krahn inequality for $p$-Hermite operators. Version arXiv:2609.01269v1 (CC BY 4.0), distributed by arXiv under the Creative Commons Attribution 4.0 International licence, http://creativecommons.org/licenses/by/4.0/, as recorded in the arXiv licence metadata for this exact version and shown on the abstract page https://arxiv.org/abs/2609.01269v1. Full licence notice: \texttt{licenses/arxiv-2609.01269-CC-BY-4.0.txt}.\par\smallskip

\noindent\textit{Editorial scope.} Counted unit: the article's proposition prop (source0.tex lines 528--531). The article's complete original proof of it is delivered with it (source0.tex lines 532--630, one \texttt{proof} environment of 99 lines). No mathematics is written, repaired or re-derived: the statement and the proof are the article's own lines, in the article's own notation, and no retained line is substituted or re-flowed.  Retained uncounted as the source-local context the proof needs: lines 499--504; lines 506--508; lines 524--526.  Nothing was omitted from any retained span, so the package carries no omission record.  No retained line contains a citation, so the wrapper carries no bibliography and no bibliography entry.  No retained line contains a figure environment, an image-inclusion command or drawing code, so no image is delivered and the package declares no asset dependency.  Editorial formatting only, in addition: an AMS \texttt{amsart} preamble that restates the article's own declarations and macros used by the retained lines, namely the environments \texttt{proposition} on separate counters, together with the standard packages those lines need.  The retained environments print in this document as Proposition 1 in order of appearance, whereas the article prints the same items with the numbering of its own full document.  The complete native source of the article as distributed in the arXiv e-print for this exact version is bundled unchanged as \texttt{sources/209091/source0.tex}.
\begin{equation}\label{3.5}
\begin{cases}
-(p-1)|w'|^{p-2}w^{''}+t|w'|^{p-2}w'=\lambda_1(\sigma)|w|^{p-2}w,\qquad t\in I_{\sigma}\\
|w'|^{p-2}w'(\sigma)+\beta|w|^{p-2}w(\sigma)=0.
\end{cases}
\end{equation}

\begin{equation}\label{3.6}
\lim\limits_{t\rightarrow -\infty}|w'|^{p-2}w'(t)e^{-\frac{t^2}{2}}=0.
\end{equation}

\begin{align}\label{def of beta}
   g(t):=\frac{|w'|^{p-1}}{w^{p-1}},
\end{align}

\begin{proposition}\label{monotone of beta prop}
    The function $g(t)$ defined in (\ref{def of beta}) is positive, strictly increasing in $(-\infty,\sigma)$,
and $g(\sigma) = \beta$.
\end{proposition}

\begin{proof} 
From equation \eqref{3.5},
a direct computation gives
\begin{align}\label{3.10}
    g'(t)&=\frac{(p-1)|w'|^{p-3}w'w''}{w^{p-1}}-(p-1)\frac{|w'|^{p-1}}{w^{p}}w'\nonumber\\
    &=\frac{1}{w^{p-1}}(\lambda_1(\sigma) w^{p-1}+t|w'|^{p-1})+(p-1)g^{\frac{p}{p-1}}\nonumber\\
    &=\lambda_1(\sigma) +tg(t)+(p-1)g(t)^{\frac{p}{p-1}}.
\end{align}
and consequently,
\begin{align}\label{g'' equation}
    g''=(t+pg(t)^{\frac{1}{p-1}})g'(t)+g(t).
\end{align}
We first show that $g(t)$ has at most one local minimum point in 
$(-\infty,\sigma]$.   Indeed, if  $t_{0}$ is a critical point of $g$, then \eqref{g'' equation} yields
\begin{align*}
    g''(t_{0})=g(t_{0})>0,
\end{align*}
so $t_{0}$ is necessarily a local minimum. Therefore,  $g$ has no local maximum points in $(-\infty,\sigma]$. If there were  two distinct local minima  $t_1< t_2$, then since $g'(t_i)=0$ and
$g''(t_i)>0$,  there would exist small 
 constants $\delta_1$ and $\delta_2$ small enough, such that $g'>0$ on  $ (t_{1}, t_{1}+\delta_1)$ and $g'<0$ on  $(t_{2}-\delta_2, t_{2})$. This would force the existence of a local  maximum  in between, a contradiction.

 Thus, only three  possibilities remain:
 {\bf Case 1:} $g'(t)>0$ for all $t\in I_\sigma$;
   {\bf Case 2:} $g'(t)\le 0$ for all $t\in I_\sigma$;
   {\bf Case 3:} There exists $t_0\in (-\infty, \sigma)$ such that
 $g'\le 0$ on $(-\infty, t_0)$ and $g'>0$ on $(t_0, \sigma)$. 
We now show that Cases 2 and 3 are impossible. Case 2 may be viewed as the limiting case of Case 3 with $t_0=\sigma$, so we treat them uniformly. Suppose $g'<0$ on $(-\infty, t_0)$, then 
$$
A:=\lim_{t\rightarrow -\infty}g(t)\in (0, +\infty].
$$
If $0<A<+\infty$,  then from \eqref{3.10} we have
$$
g'(t)\to -\infty, \quad \text{as $t \to -\infty$},
$$
which implies $g(t)\to +\infty$, contradicting with $0<A<+\infty$.
% using \eqref{3.6} and $w'(t)<0$, we have
% $$\lim_{t\rightarrow-\infty}w(t)^{p-1}e^{-\frac{t^{2}}{2}}=\lim_{t\rightarrow-\infty}\frac{1}{g(t)}|w'(t)|^{p-1}e^{-t^2/2}=0.$$ 
% Thus,
% \begin{align*}
% \lim_{t \to -\infty}g(t)
% &=\lim_{t \to -\infty} \frac{|w'(t)|^{p-1} e^{-\frac{t^2}{2}}}{|w(t)|^{p-1} e^{-\frac{t^2}{2}}}\\
% &= \lim_{t \to -\infty} \frac{(p-1)|w'|^{p-3} w' w'' - t |w'|^{p-1}}{(p-1)|w|^{p-2} w' - |w|^{p-1} t} \\
% &= \lim_{t \to -\infty} \frac{-(p-1)|w'|^{p-2} w'' - t |w'|^{p-1}}{(p-1)|w|^{p-2} w' - |w|^{p-1} t} \\
% &= \lim_{t \to -\infty} \frac{\lambda_1(\sigma) |w|^{p-1}}{(p-1)|w|^{p-2} w' - |w|^{p-1} t} \\
% &= \lim_{t \to -\infty} \frac{\lambda_1(\sigma)}{-(p-1)g^{1/p-1}- t} = 0,
% \end{align*}
% which contradicts $A>0$. 
Hence $A=+\infty$.  Since
\begin{align*}
    \lambda_1(\sigma)& = \frac{\int_{-\infty}^{\sigma} |w'|^p e^{-\frac{t^2}{2}} \, dt + \beta |w(\sigma)|^pe^{-\frac{\sigma^2}{2}}}{\int_{-\infty}^{\sigma} w^p e^{\frac{-t^2}{2}} \, dt}\\
&= \frac{\int_{-\infty}^{t} |w'|^p e^{-\frac{t^2}{2}} \, dt + \int_{t}^{\sigma} |w'|^p e^{-\frac{t^2}{2}} \, dt + \beta |w(\sigma)|^p e^{-\frac{\sigma^2}{2}} }{\int_{-\infty}^{t} w^p e^{\frac{-t^2}{2}}\, dt + \int_{t}^{\sigma} w^p e^{\frac{-t^2}{2}}\, dt}
\end{align*}
 and for any $t\in I_\sigma$,
\begin{align*}
    \lim_{t\rightarrow-\infty}\frac{ \int_{t}^{\sigma} (w')^p e^{-\frac{t^2}{2}} \, dt + \beta w^p(\sigma) e^{-\frac{\sigma^2}{2}} }{\int_{t}^{\sigma} w^p e^{\frac{-t^2}{2}}\, dt}=\lambda_1(\sigma),
\end{align*}
there exists $t_{3}<\sigma$ such that for any $t<t_{3}$,
\begin{align}
    \frac{ \int_{t}^{\sigma} |w'|^p e^{-\frac{t^2}{2}} \, dt + \beta |w(\sigma)|^p e^{-\frac{\sigma^2}{2}} }{\int_{t}^{\sigma} w^p e^{\frac{-t^2}{2}}\, dt}<\lambda_1(\sigma)+1.
\end{align}
Now choose a positive constant  $K>0$ such that $K>\lambda_1(\sigma)+1$. Since  $A=+\infty$, there exists    $t_{4}$ such that 
\begin{align*}
        g(t)>K^{\frac{p-1}{p}} \quad \text{for any $t\leq t_{4}$},
\end{align*}
 hence
 $|w'(t)|>K^{\frac{1}{p}}w(t)$. Therefore, for $t\le t_4$
\begin{align*}
    \frac{\int_{-\infty}^{t} |w'|^p e^{-\frac{t^2}{2}} \, dt}{\int_{-\infty}^{t} w^p e^{-\frac{t^2}{2}} \,dt} > K.
\end{align*}
Let  $t_5=\min\{t_3,t_4\}$ and define
\begin{align*}
    a_{1}(t):=&\int_{t}^{\sigma} |w'|^p e^{-\frac{t^2}{2}} \, dt + \beta |w(\sigma)|^p e^{-\frac{\sigma^2}{2}},\\
     a_{2}(t):=&\int_{-\infty}^{t} |w'|^p e^{-\frac{t^2}{2}} \, dt,\\
     b_{1}(t):=&\int_{t}^{\sigma} w^p e^{\frac{-t^2}{2}}\, dt,\\
     b_2(t):=&\int_{-\infty}^{t} w^p e^{-\frac{t^2}{2}} \, dt.
\end{align*}
Then for  $t\leq t_5$,
\begin{align*}
   \frac{a_2(t)}{b_2(t)} > K, \quad \ell (t):= \frac{a_1(t)}{b_1(t)}<\lambda_1(\sigma)+1<K. 
\end{align*}
Therefore we conclude that
\begin{align*}
\lambda_1(\sigma)=\frac{a_1 + a_2}{b_1 + b_2} > \frac{K b_2 + \ell b_1}{b_2 + b_1} > \ell = \frac{a_1}{b_1}=\frac{\int_{t}^{\sigma} |w'|^p e^{-\frac{t^2}{2}} \, dt + \beta |w(\sigma)|^p e^{-\frac{\sigma^2}{2}}}{\int_{t}^{\sigma} w^p e^{\frac{-t^2}{2}}\, dt}.
\end{align*}
Choosing a testing function 
\begin{align*}
    \widetilde{w}(t) = 
\begin{cases} 
w(t), & t \in (t_5, \sigma), \\
w(t_5), & t \in (-\infty, t_5)
\end{cases}
\end{align*}
for $\l_1(\sigma)$, we have 
\begin{align*}
    \lambda_1(\sigma)\leq \frac{\int_{t_5}^{\sigma} |w'|^p e^{-\frac{t^2}{2}} \, dt + \beta |w(\sigma)|^p e^{-\frac{\sigma^2}{2}}}{\int_{t_5}^{\sigma} w^p e^{\frac{-t^2}{2}}\, dt+w(t_{5})^p\int_{-\infty}^{t_{5}}e^{-\frac{t^{2}}{2}}dt}< \frac{\int_{t_5}^{\sigma} |w'|^p e^{-\frac{t^2}{2}} \, dt + \beta |w(\sigma)|^p e^{-\frac{\sigma^2}{2}}}{\int_{t_5}^{\sigma} w^p e^{\frac{-t^2}{2}}\, dt}<\lambda_1(\sigma),
\end{align*}
which is a contradiction. This rules out Cases 2 and 3, so Case 1 must hold, i.e.,
 $g'(t)>0$.
\end{proof}

\end{document}
